\documentclass[a4paper, 12pt]{article}

\usepackage[top=2cm, bottom=2cm, left=2.5cm, right=2.5cm]{geometry}
\usepackage[utf8]{inputenc}
\usepackage{amsmath, amsfonts, amssymb}
\usepackage{graphicx} % inserir figuras - \includegraphics[scale=•]{•}
\usepackage{float} % ignorar regras de tipografia e inserir figura aonde queremos.
\usepackage[brazil]{babel} % Trocar Figure para Figura.
\usepackage{indentfirst}
\pagestyle{empty}


\begin{document}
\begin{figure}[H]
	\includegraphics[scale=0.9]{UnB_CiC_Logo.jpg}
\end{figure}
\noindent\rule{\textwidth}{0.4pt}
\begin{center}
	\textbf{{\Large Introdução à Ciência da Computação - 113913}} \newline \newline
	\textbf{{\large Prova 1} \\
	\vspace{9pt}
	{\large Questão A}} \\
	\noindent\rule{\textwidth}{0.4pt}
	\newline
\end{center}

\textbf{{\large Observações:}}
\begin{itemize}
	\item As provas serão corrigidas por um corretor automático, portanto é necessário que as entradas e saídas do seu programa estejam conforme o padrão especificado em cada questão (exemplo de entrada e saída).
	\item Por exemplo, não use mensagens escritas durante o desenvolvimento do seu código como “Informe a primeira entrada".
	\item Estas mensagens não são tratadas pelo corretor, portanto a correção irá resultar em resposta errada, mesmo que seu código esteja correto.
	\item Serão testadas várias entradas além das que foram dadas como exemplo, assim como as listas.
	\item Assim como as listas, as provas devem ser feitas na versão Python 3 ou superior.
	\item Cada questão (A e B) vale 50\% da nota da Prova 1.
	\item Leia com atenção e faça \textbf{exatamente} o que está sendo pedido.


\end{itemize}
\newpage % Questão A 
\begin{center}
\textbf{{\Large Questão A - Saneamento em bacias hidrográficas}}
\end{center}

\vspace{5pt} 
Bacia Hidrográfica é a área ou região de drenagem de um rio principal e seus afluentes. É a porção do espaço em que as águas das chuvas, das montanhas, subterrâneas ou de outros rios escoam em direção a um determinado curso d’água, abastecendo-o. Os serviços de saneamento devem ser analisados pela perspectiva da Bacia Hidrográfica (BH).  Muitos municípios brasileiros não possuem tratamento de esgotos adequado ou sequer disponibilizam o serviço para sua população.O lançamento desses efluentes nos corpos hídricos comprometem a qualidade e os usos das águas, causando implicações danosas à saúde pública e ao equilíbrio do meio ambiente.\newline
Ache o município que apresenta a maior vazão de esgoto sem tratamento, o município com a maior taxa de tratamento de esgoto, o município com maior percentual da população urbana atendida com esgoto sanitário e o valor médio que os municípios terão que investir em coleta e tratamento de esgoto nos próximos 20 anos para atenuar o problema.
\newline \newline
\textbf{{\large Entrada}} \newline
Um número \textbf{N} que é o número de municípios de uma BH. \newline
\textbf{N} linhas, cada linha com os seguintes dados do município:\newline Nome, população urbana, vazão de efluente sem coleta e sem tratamento em L/s, vazão de efluente com coleta e com tratamento em L/s, a população atendida por sistema de coleta e tratamento de esgoto e o quanto cada município deve investir nos próximos 20 anos para tentar resolver o problema.
\newline \newline
\textbf{{\large Saída}} \newline
Nome do município que apresenta a maior vazão de esgoto sem tratamento.\newline
Nome do município com a maior taxa de tratamento de esgoto. \newline
O município com maior percentual da população urbana atendida com esgoto sanitário. \newline
O valor médio que os municípios investem em coleta e tratamento de esgoto com 2 casas decimais.
\newline \newline
\newline
\begin{table}[H]
\centering
\begin{tabular}{|l|l|}
\hline
\textbf{Entrada}                                                                                                              & \textbf{Saída}                                                      \\ \hline
\begin{tabular}[c]{@{}l@{}}4\\ ES 130 35.5 20.5 80 1000\\ TU 150 100.0 98.0 100 3900\\ DO 800 290.0 100.0 100 8000\\ MT 350 190.0 80.0 330 2500 \end{tabular}                                                                                                          & \begin{tabular}[c]{@{}l@{}}DO\\ TU\\ MTO\\ R\$ 3850.00\end{tabular}        \\ \hline
\begin{tabular}[c]{@{}l@{}}4\\ Ariquemes-RO 85770 125.7 2.8 2000 144502208.57\\ Cabixi-RO 2771 4.3 0.0 0 5550611.37\\ Ji-Paraná-RO 115123 145.0 8.0 6351 162456203.88\\ Cajubim-RO 13520 19.5 0.2 420 25171237.94\end{tabular}                                            & \begin{tabular}[c]{@{}l@{}}Ji-Paraná-RO\\ Ji-Paraná-RO\\ Ji-Paraná-RO\\ R\$ 84420065.44\end{tabular} \\ \hline

\end{tabular}
\caption{Questão A}
\end{table}
\flushright
\textbf{\Large Boa Prova!}
\end{document}